\subsection{Funciones predictivas y entradas recurrentes}
Para historia preparada $ v $, persistencia predice $ v_L $. La media móvil usa $ k^{-1}\sum_{i=L-k+1}^L v_i $ y suavizamiento exponencial directo el último estado de la recurrencia actualizada por disponibilidad. La regresión ridge \cite{hoerlkennard} minimiza $\sum_j(x_j-w^\top v_j-b)^2+\lambda\|w\|_2^2$ sobre ejemplos fuente preparados $ j $. NLinear predice $ v_L+w^\top(v-v_L\mathbf1)+b $. Estas expresiones operan en escala normalizada. La misma transformación inversa precede las métricas. Su sencillez ofrece referencias directas para valorar cuánto aporta una representación más rica.

Sea $ q_i $ el triplete de valor, indicador y edad por servicio al paso $ i $. La GRU de grafo consume $[q_i,D_{\rm out}^{-1}Gq_i,D_{\rm in}^{-1}G^\top q_i]$, con filas cero cuando no hay vecinos. La GRU local usa identidades en ambas posiciones. La función recurrente es $ h_i=\operatorname{GRU}(u_i,h_{i-1})$ y la cabeza final $ w^\top h_L+b $. En GRU-D, los decaimientos son $ g_x=\exp[-\operatorname{ReLU}(W_x\delta+b_x)]$ y $ g_h=\exp[-\operatorname{ReLU}(W_h\delta+b_h)]$. La celda recibe $[m_iv_i+(1-m_i)g_xv_i,m_i]$ y estado previo $ g_h\odot h_{i-1}$, iniciado en cero. Se atenúa el valor preparado hacia cero normalizado. No se recupera por ello la medición oculta. Los agregados de vecindad y su disponibilidad temporal se distinguen de los parámetros entrenables.

PatchTST aplica normalización reversible de instancia \cite{revin}, proyecta parches de longitud cuatro con paso fijo e información posicional, y asigna los parches codificados a un objetivo. Cada cabeza usa $\operatorname{softmax}(QK^\top/\sqrt{d_k})V $ sobre las representaciones. La proyección final y normalización inversa devuelven el extremo a la escala de la entrada preparada. La descripción corresponde a la receta de un canal ejecutada. BiTGraph y GinAR históricos conservan adaptaciones de código de autores y controles de validación finita. Sus resultados documentan esas implementaciones locales.

Si una ventana no contiene mediciones nuevas, su último valor, su media y el nivel recursivo se definen mediante arrastre o respaldo causal. Los modelos con descriptores también pueden responder a procedencia y edad. Estas ventanas participan en la misma evaluación emparejada que las más completas. El flujo preserva sus objetivos originales y no elimina selectivamente errores grandes. Se compara así precisión puntual condicionada a las máscaras, sin evaluar detección y recuperación de averías reales.
